%  LaTeX support: latex@mdpi.com 
%  For support, please attach all files needed for compiling as well as the log file, and specify your operating system, LaTeX version, and LaTeX editor.

%=================================================================
\documentclass[sustainability,article,submit,pdftex,moreauthors]{Definitions/mdpi} 

%=================================================================
% MDPI internal commands
\firstpage{1} 
\makeatletter 
\setcounter{page}{\@firstpage} 
\makeatother
\pubvolume{1}
\issuenum{1}
\articlenumber{0}
\pubyear{2022}
\copyrightyear{2022}
%\externaleditor{Academic Editor: Firstname Lastname}
\datereceived{} 
\dateaccepted{} 
\datepublished{} 
%\datecorrected{} % Corrected papers include a "Corrected: XXX" date in the original paper.
%\dateretracted{} % Corrected papers include a "Retracted: XXX" date in the original paper.
\hreflink{https://doi.org/} % If needed use \linebreak
\graphicspath{{figures/}} % path to folder with figures
%\doinum{}

\usepackage{listings}
\usepackage{xcolor}

\definecolor{codegreen}{rgb}{0,0.6,0}
\definecolor{codegray}{rgb}{0.5,0.5,0.5}
\definecolor{codepurple}{RGB}{204, 0, 153}
\definecolor{backcolour}{RGB}{255, 250, 230}
%    
\lstdefinestyle{mystyle}{
    backgroundcolor=\color{backcolour}, 
    commentstyle=\color{codegreen},
    keywordstyle=\color{magenta},
    numberstyle=\tiny\color{codegray},
    stringstyle=\color{codepurple},
    basicstyle=\ttfamily\footnotesize,
    breakatwhitespace=false,         
    breaklines=true,                 
    captionpos=b,                    
    keepspaces=true,                 
    numbers=left,                    
    numbersep=5pt,                  
    showspaces=false,                
    showstringspaces=false,
    showtabs=false,                  
    tabsize=2
}

\lstset{style=mystyle}
	
\usepackage{soul} % highlighting text
%------------------------------------------------------------------
% The following line should be uncommented if the LaTeX file is uploaded to arXiv.org
%\pdfoutput=1

%=================================================================
% Add packages and commands here. The following packages are loaded in our class file: fontenc, inputenc, calc, indentfirst, fancyhdr, graphicx, epstopdf, lastpage, ifthen, lineno, float, amsmath, setspace, enumitem, mathpazo, booktabs, titlesec, etoolbox, tabto, xcolor, soul, multirow, microtype, tikz, totcount, changepage, attrib, upgreek, cleveref, amsthm, hyphenat, natbib, hyperref, footmisc, url, geometry, newfloat, caption

%=================================================================
%% Please use the following mathematics environments: Theorem, Lemma, Corollary, Proposition, Characterization, Property, Problem, Example, ExamplesandDefinitions, Hypothesis, Remark, Definition, Notation, Assumption
%% For proofs, please use the proof environment (the amsthm package is loaded by the MDPI class).

%=================================================================
% Full title of the paper (Capitalized)
\Title{Mapping gravity anomalies and geologic heterogeneity of Venezuela linked to seismotectonics of shallow-focus earthquakes along the plate margins}

% MDPI internal command: Title for citation in the left column
\TitleCitation{Mapping gravity anomalies and geologic heterogeneity of Venezuela linked to seismotectonics of shallow-focus earthquakes along the plate margins}

% Author Orchid ID: enter ID or remove command
\newcommand{\orcidauthorA}{0000-0002-5759-1089} % Add \orcidA{} behind the author's name
\newcommand{\orcidauthorB}{0000-0002-6461-1551} % Add \orcidB{} behind the author's name

% Authors, for the paper (add full first names)
\Author{Polina Lemenkova $^{1}$*\orcidA{} and Olivier Debeir $^{1}\orcidB{}$}

% MDPI internal command: Authors, for metadata in PDF
\AuthorNames{Poina Lemenkova and Olivier Debeir}

% MDPI internal command: Authors, for citation in the left column
\AuthorCitation{Lemenkova, P.; Debeir, O.}
% If this is a Chicago style journal: Lastname, Firstname, Firstname Lastname, and Firstname Lastname.

% Affiliations / Addresses (Add [1] after \address if there is only one affiliation.)
\address{%
$^{1}$ \quad Université Libre de Bruxelles (ULB), École polytechnique de Bruxelles (Brussels Faculty of Engineering), Laboratory of Image Synthesis and Analysis (LISA). Building L, Campus de Solbosch, ULB -- LISA CP165/57, Avenue Franklin D. Roosevelt 50, Brussels 1000, Belgium.}

% Contact information of the corresponding author
\corres{Correspondence: polina.lemenkova@ulb.be; Tel.: +32471860459 (P.L.)}

% The commands \thirdnote{} till \eighthnote{} are available for further notes

%\simplesumm{} % Simple summary

% Abstract (Do not insert blank lines, i.e. \\) 
\abstract{This paper presents a cartographic framework based on algorithms of GMT codes for visualisation of spatial data. The aim is to map a seismically endangered areas in selected countries of South America. Spatial data were integrated in the project including GEBCO, EGM-2008, IRIS, USGS geological layers. The extent of the countries was extracted and data were processed iteratively in the console of GMT, converted by GDAL, formatted, and mapped for geophysical data visualisation. The GIS methods were applied to visualise, map and analyse layers through a spatial analysis and cartographic overlay techniques. The 2,000 samples of earthquake events were obtained from the IRIS seismic database covering the 25-year time span (1997-2021). The data were converted to ‘ngdc’ format, visualized by 'psxy' module and assessed for seismicity. The approach to linking geological, topographic and geophysical data using GMT algorithms aimed to map correlations between the geophysical phenomena, tectonic processes, geological objects, seismicity and earthquakes. Andes is one the of the most seismically active region of the world, therefore, practical application of GMT scripts consists in contribution to the geological risk assessment and hazard prevention in real time regime. In response to the effects of topographic variability and tectonic setting, the distribution of earthquakes indicated the closeness to the margins of the colliding South America and Caribbean tectonic plates, mountain geomorphology and geophysics with variations in depth and magnitude of seismic events. Geologically, the result showed that the most important factors that have a significant influence on the seismicity in Venezuela, Trinidad and Tobago are plate subduction and associated complex geological processes, geodynamic heterogeneity of the upper mantle and geomorphology, the extent of the Tobago Trough, Cordillera de Mérida, segments of Andean orogeny and Lara-Falcón Basin. The proposed method integrative the algorithms of the GMT scripts with GIS considers multiple cartographic techniques as effective tool for mapping spatial datasets and rapid data processing in iterative regime. In this context, using GMT and GIS for finding similarity between the regional earthquake distribution and geological and topographic setting is essential for hazard risk assessment. This study can serve as a basis for predictive seismic analysis in geological vulnerable regions, such as Venezuela. Besides technical demonstration of GMT algorithms, the study also contributes to the geologic risk assessment, geophysical mapping and seismic hazard assessment in South America. Full algorithms of GMT scripts are available in the author’s GitHub repository.}

\keyword{Risk, hazard, sustainability, earthquake, seismicity, cartography, GMT, geophysics} 

\PACS{91.10.Da, 91.10.-v, 91.10.Jf, 91.10.Op, 91.30.Dk, 93.85.Pq, 91.70.-c}
\MSC{86Axx, 86A04, 86A15, 86A30, 86A60, 86-XX, 86-04, 86-08}
\JEL{Y91, Q00, Q01, Q2, Q20, Q24, Q3, Q35, Q5, Q50, Q51, Q54, Q55, Q56, C6, C61, C63}

%%%%%%%%%%%%%%%%%%%%%%%%%%%%%%%%%%%%%%%%%
\begin{document}

\section{Introduction}

\subsection{Background and motivation}

Nowadays, geologic hazards are acknowledged as one of the most disastrous and risk events for cities locate din the mountainous regions. These involve a series of negative consequences for nature (landslides, tsunami, rock falls, land mass movements) and complex social problems (destruction of houses and infrastructure, injures and losses of human lives) caused by seismic risks \citep{Laffaille,Schmitz2020,Weber}. This study focuses on the Venezuela, and the Republic of Trinidad and Tobago, South America with multiple cities at risk of seismic hazards (Fig. 1). Early prediction of seismic risks in geologically vulnerable regions has become more valuable in a large variety of time-critical geospatial applications focused on evaluating geological hazards aimed to mitigate earthquake hazards at the urban scale. 


\begin{figure}[H]\centering
	\includegraphics[width=12.0 cm]{Fig_01}
	\caption{Topographic map of the study area comprising the two countries of South America: the Venezuela and the Trinidad and Tobago. Mapping: GMT. Source: author. \label{fig1}}
\end{figure}  

To build an effective geological mapping and modelling for seismic hazard prediction, geophysical setting, processes and activities can be characterised by a spatio-temporal composition of the constituent geophysical and geological forces, evolutional processes and their complex interactions. Such linkages results in the actual location and distribution of the seismic sources of earthquakes with regard to the location of cities, and information on these processes can be effectively used for seismic prognosis and earthquake modelling. In recent decades, research demonstrated that geological and geophysical mapping creates a basis for natural hazard risk assessment. The data derived from geophysical and geological surveys can be used as a background information and recognition of complex geological actions and processes. 

Advances in cartographic development and mapping made an important application of spatial data analysis become real: predicting and prognosis of seismic activities or imminent earthquake events from the observed events in the past. Thus, using the retrospective data we can make the prognosis of possible seismicity and help avoiding the risk for people at risk due to their living locations close ti the seismic sources. Many intelligent systems aimed at geological prognosis can benefit from real-time mapping application and updated cartographic datasets used for seismic prediction. For instance, a systematic examination of the correlations and links among the various geophysical and geological determinants underlying disasters may be used to prevent socioeconomic risks in the geologically unstable zones. In the earthquake analysis, the capability of predicting the possible events, their magnitude and focal depth is highly desirable for real-time risk assessment. By integration of geological data with socio-economic analysis can be used for effective management of possible risks, preventing damages and dangers. 

Nevertheless, though geophysical mapping is a very important task of such seismically active regions as South America, it is a quite new topic for scripting cartography. To the best of our knowledge, the majority of similar studies use the conventional GIS rather than scripts for data visualisation. Using GIS enables to plot maps in the traditional regime which requires handmade routine and cannot be acceptable for rapid real-time mapping. Using GIS limits the types of cartographic activities that can be used for rapid geological mapping and seismic prediction as well as how the data from various sources can be handled.  In contrast, running scripts presents a rapid and repeatable procedure of mapping that can be integrated for real-time risk assessment and hazard analysis. Thus, the more flexible cartographic data processing, enabled by scripts, is more desirable and valuable for risk assessment if it focuses on retrospective multi-source geological datasets. 

Early detection of seismic risks can be solved by evaluating and visualising different geophysical determinants that affect disaster fatality at the country level, namely seismic hazard strength, population data including distribution of locations and population density, demographics and socioeconomics data as well as geological background layers. The approach presented in this study aims to solve the problem of traditional mapping by introducing the integrative study on scripts developed in the Generic Mapping Tools (GMT) scripting toolset. Such a framework, which is different from the traditional GIS approach, presents a quick mapping from the console using the command line and codes written by the GMT syntax. 

Specifically, this paper is focused on the two seismically active countries of the South America: 1) Venezuela; 2) Trinidad and Tobago. The cartographic aim is to demonstrate the use of GMT scripts in a combination with GIS techniques. The geological aim is to understand the correlation between the geological setting, topography and geophysical anomalies to the regional seismicity by analysis of factors that determine the distribution, magnitude and depth of earthquakes. The study applied the GMT scripting cartographic techniques and QGIS to visualize the high-resolution open sources raster grids and vector data. The semantic key of this approach is to utilize the variety of datasets to build the project with observed geological and geophysical processes and structures as context to predict the next possible seismic events, the intension or the possible effect from the earthquakes on the population of both countries: Venezuela and the Republic of Trinidad and Tobago. 

\subsection{Study area}

The north-western and north-eastern regions of the study area are vulnerable to earthquakes due to the tectonic activity of the lithosphere plates \citep{Barr,Hayes}. In particular, the Caribbean Sea region is a place with complex geologic setting (Fig. 2) where several tectonic plates experience collision, interaction and subduction: Caribbean, Cocos, Nazca and South America plates. This results in high seismic activity in the region under the influence of tectonic processes which his reflected in current topography, geology and geomorphology of South America in general and Venezuela in particular, as has been discussed in previously published relevant papers \citep{Diebold,Lemenkova2019b,Fajardo,Lemenkova2020b,Perez1981}. 

\begin{figure}[H]\centering
	\includegraphics[width=14.0 cm]{Fig_02}
	\caption{Geologic provinces in Venezuela, Trinidad and Tobago and neighbouring countries. Mapping: QGIS. Source: author.\label{fig2}}
\end{figure}  

Venezuela is a country located in the north of South American continent where a collision of South American and Caribbean tectonic plates takes place causing complex regional geologic pattern (Fig. 3). The geodynamics of the interacted plates results in the formation of thrust belts and foreland basins in northern Venezuela that are formed by the interplay of geologic layers. Thus, the outcrops of Precambrian (pC), Quaternary (Q), Triassic (T), Paleozoic-Mesozoic intrusives (PMi) are dominating in the central regions of the country while regional outcrops of Cretaceous (K), Paleozoic metamorphics (PZm) and Mesozoic-Cenozoic intrusives (MCi) are in the Falcón Basin (Fig. 3). 

The repetitive occurrence of seismic events, such as volcanism and earthquakes in north-western regions of Venezuela and Trinidad and Tobago are recorded as series of experienced events throughout the recent 25 years by IRIS database. The ongoing effects of earthquakes in Venezuela \citep{Audemard2006,Baumbach,Colon,Pousse2018,Rial,Yamazaki} and coastal regions of Trinidad and Tobago \citep{Deville,Latchman}, made geophysical mapping important task aimed at seismic hazard mapping, earthquake prevention and mitigation, and risk assessment \cite{Bucci}.

Tectonic active movements cause the formation of the orogenic mountain belt and topography, bathymetry complicated with deep-sea trenches and troughs in the margins of the Pacific Ocean and the Caribbean Sea along the subduction path \citep{Gough,Lemenkova2020a}. Complex geological history is especially notable in the NW Venezuela, Falcón Basin, where the interactions between the lithosphere plates and subducting Caribbean slab cause the associated seismic activity and geologic processes: outcrop of igneous intrusive bodies, crustal thinning, decreased Moho depth, gravity anomalies, rifting \cite{Bezada}. \cite{Mazuera} identified crustal thinning beneath the Falcón Basin along the western extension of the Oca-Ancón fault system which they interpreted as a back-arc basin, and suture zones between the Proterozoic and Paleozoic provinces (Ouachita-Marathon-related suture), and Paleozoic and Meso-Cenozoic terranes (peri-Caribbean suture) where seismic velocity changed. 

\begin{figure}[H]\centering
	\includegraphics[width=14.0 cm]{Fig_03}
	\caption{Geologic map of Venezuela, Trinidad and Tobago and neighboring countries. Mapping: QGIS. Source: author.\label{fig3}}
\end{figure} 

\subsection{Actuality and objectives}

The actuality of the undertaken study can be illustrated by the social vulnerability towards geologic risks. Venezuela is one of the most vulnerable countries to earthquakes compared to other countries in Amazon Basin. It is reported that since the 19th century the tsunami in the Caribbean Sea resulted in several thousands of lives lost and many people exposed at risk and damage from the landslides \cite{Harbitz}. The frequency and magnitude of earthquakes are associated with geologic and geophysical setting and tectonic interaction of the lithosphere plates that affect seismic activity and tectonic-related hazards that affect the lives of people and infrastructure \citep{Masy,Perez2018,Suarez} Masy et al. 2011. High tsunami hazard and risk of associated earthquakes and submarine volcanos exist in the Caribbean region and the Venezuela. 

Thus, the Caribbean-South American plate boundary is a very wide tectonically active zone \cite{Audemard1993,Backe} with active seismicity at various depths. For instance, anomalous earthquake series took place in the east coast of the Trinidad with reported depth of the main shock, 53 km \cite{Marshall}. Seismicity deeper than 40 km is reported in the southern end of the Lesser Antilles subduction zone under Trinidad and Tobago and the active island arc of the Lesser Antilles \cite{Reinoza}. 

Therefore, the integrated mapping using various multi-sources data is motivated by complex interplay among geological, topographic, tectonic and geophysical setting \cite{Lemenkova2021a}. For instance, tectonic architecture with the formation of deep basins and erosional troughs affects submarine geomorphology \cite{Lemenkova2021c} which can be interpreted using geophysical and gravimetric measurements \cite{Lemenkova2021e}. The earthquake-related aftershocks, hazards and stresses, such as landslides affecting geotechnical infrastructure and bridges \cite{Salcedo}, demolished buildings and lost of human lives \cite{Rodriguez2018}, aseismic slips and creeps \cite{Pousse2016} affect both nature and social sector of Venezuela and require regular monitoring,  accurate mapping and precise risk and vulnerability assessment. 

Previous cartographic investigations have been carried out using traditional GIS applications in geologic and environmental studies \citep{Gonzalez,Klauco2013,Klauco2017,Lemenkova2020d,Hackley,Schmitz2021,Sun}. However, none of them have focused on the multi-disciplinary application combining geologic research with machine learning methods of mapping Venezuela. 

Compared to the traditional GIS application often used in previous studies, using scripting techniques in geological mapping has not been adequately presented and reported in the existing literature so far. At the same time, scripting techniques employed for visualization of the complex geophysical and geologic phenomena enable to highlight correlations through objective machine-based graphics and automated visualization of the determinants and factors \citep{Lemenkov2021a,Lemenkov2021b}. With this aim, the Generic Mapping Tools (GMT) approach of scripting cartographic toolset combined with QGIS and its plugins affecting is applied for geophysical mapping of Venezuela, Trinidad and Tobago to show relative distribution of the seismicity with regards to the geophysical, topographic and geologic phenomena. 

%%%%%%%%%%%%%%%%%%%%%%%%%%%%%%%%%%%%%%%%%%
\section{Materials and Methods}

The study area encompasses the following spatial extent: 59,5°-74°W, 0°-12.5°N (Fig. 1). Topographic, geophysical, seismic and geologic sources of data were used in this study presenting a multi-disciplinary cartographic analysis. The topographic data were collected through the General Bathymetric Chart of the Oceans (GEBCO) open access repository with NetCDF raster grid \cite{GEBCO} which presents the most comprehensive data for topography and bathymetry of the Earth \cite{Schenke}. GEBCO data are widely used in geosciences due to its unprecedented resolution and precision \citep{Mcmichael,Lemenkova2020c,Lemenkova2022c,Cortina}. 

The geological data (Fig. 2 and 3) have been plotted using the QGIS, because this open source software is compatible with the ArcGIS data format that have been used in geologic maps as a .shp vector layers. The visualization of these maps followed the traditional mapping process of the QGIS \cite{Lemenkova2022a}. In the first stage, vector layers of the geologic provinces and outcrops were visualized with in a Graphical User Interface (GUI). In the second stage, the study area was designated using the overlay of the Digital Chart of the World (DCW), which is a comprehensive 1:1,000,000 scale vector base map uploaded into the QGIS project. 

Maps of topographic, seismic and geophysical data (Figs. 1, 4–7) were made using scripting technique of the GMT \cite{Wessel} following the existing examples of scripting \citep{Lemenkova2019a,Lemenkova2022b}. The methodological workflow included 5 algorithms by GMT codes as follows: 

\begin{lstlisting}[language=bash, caption=\textbf{Algorithm of GMT code used for plotting Fig. \ref{fig1} (topographic map)}]
#!/bin/sh
# Purpose: shaded relief grid raster map from the GEBCO dataset (here: Venezuela)
# GMT modules: gmtset, gmtdefaults, grdcut, makecpt, grdimage, psscale, grdcontour, psbasemap, gmtlogo, psconvert
# http://soliton.vm.bytemark.co.uk/pub/cpt-city/arendal/tn/arctic.png.index.html
# GMT set up
gmt set FORMAT_GEO_MAP=dddF \
# Overwrite defaults of GMT
gmtdefaults -D > .gmtdefaults
#chsh -s /bin/bash
chsh -s /bin/zsh
gmt grdcut GEBCO_2019.nc -R286/300.5/0/12.5 -Gve_relief.nc
gmt grdcut ETOPO1_Ice_g_gmt4.grd -R286/300.5/0/12.5 -Gve_relief1.nc
gdalinfo ve_relief.nc -stats
# Minimum=-4947.963, Maximum=5535.625, Mean=162.220, StdDev=788.097
# create mask of vector layer from the DCW of country's polygon
gmt pscoast -R286/300.5/0/12.5 -Dh -M -EVE > ve.txt
# Make color palette
gmt makecpt -Carctic.cpt > pauline.cpt
# Generate a file
ps=Topography_VE.ps
# Make background transparent image
gmt grdimage ve_relief.nc -Cpauline.cpt -R286/300.5/0/12.5 -JM6i -P -I+a15+ne0.75 -t20 -Xc -K > $ps
# Add isolines
gmt grdcontour ve_relief1.nc -R -J -C500 -W0.1p -O -K >> $ps
# Add coastlines, borders, rivers
gmt pscoast -R -J -P \
    -Ia/thinner,blue -Na -N1/thickest,darkred -W0.1p -Df -O -K >> $ps
# CLIPPING
# 1. Start: clip the map by mask to only include country
gmt psclip -R286/300.5/0/12.5 -JM6.0i ve.txt -O -K >> $ps
# 2. create map within mask
# Add raster image
gmt grdimage ve_relief.nc -Cpauline.cpt -R286/300.5/0/12.5 -JM6.0i -I+a15+ne0.75 -Xc -P -O -K >> $ps
# Add isolines
gmt grdcontour ve_relief1.nc -R -J -C1000 -Wthinner,darkbrown -O -K >> $ps
# Add coastlines, borders, rivers
gmt pscoast -R -J \
    -Ia/thinner,blue -Na -N1/thicker,tomato -W0.1p -Df -O -K >> $ps
# 3: Undo the clipping
gmt psclip -C -O -K >> $ps
# Add color barlegend
gmt psscale -Dg286/-1.0+w15.2c/0.4c+h+o0.0/0i+ml -R -J -Cpauline.cpt \
    -Bg500a1000f100+l"Color scale 'arctic': global bathymetry/topography relief [-5000 to 4000, mixed, RGB, 110 segments]" \
    -I0.2 -By+lm -O -K >> $ps    
# Add grid
gmt psbasemap -R -J \
    --MAP_FRAME_AXES=WEsN \
    --FORMAT_GEO_MAP=ddd:mm:ssF \
    -Bpx4f2a2 -Bpyg4f2a2 -Bsxg4 -Bsyg2 \
    -B+t"Topographic map of Venezuela, Trinidad and Tobago" -O -K >> $ps    
# Add scalebar, directional rose
gmt psbasemap -R -J \
    -Lx12.7c/-2.2c+c50+w200k+l"Mercator projection. Scale: km"+f \
    -UBL/-5p/-65p -O -K >> $ps
# Texts
gmt pstext -R -J -N -O -K \
-F+jTL+f10p,21,darkred+jLB+a-0 -Gwhite@60 >> $ps << EOF
290.1 11.8 Peninsula
EOF
# insert map
# Countries codes: ISO 3166-1 alpha-2. Continent codes AF (Africa), AN (Antarctica), AS (Asia), EU (Europe), OC (Oceania), NA (North America), or SA (South America). -EEU+ggrey
gmt psbasemap -R -J -O -K -DjBL+w3.3c+stmp >> $ps
read x0 y0 w h < tmp
gmt pscoast --MAP_GRID_PEN_PRIMARY=thinner,white -Rg -JG293/6N/$w -Da -Gbrown -A5000 -Bg -Wfaint -ESA+gpeachpuff -EVE+gyellow -Slightsteelblue3 -O -K -X$x0 -Y$y0 >> $ps
#gmt pscoast -Rg -JG12/5N/$w -Da -Gbrown -A5000 -Bg -Wfaint -ECM+gbisque -O -K -X$x0 -Y$y0 >> $ps
gmt psxy -R -J -O -K -T  -X-${x0} -Y-${y0} >> $ps
# Add GMT logo
gmt logo -Dx6.2/-2.9+o0.1i/0.1i+w2c -O -K >> $ps
# Add subtitle
gmt pstext -R0/10/0/15 -JX10/10 -X0.5c -Y5.0c -N -O \
    -F+f10p,Helvetica,black+jLB >> $ps << EOF
2.3 13.3 SRTM/GEBCO 15 arc sec resolution global terrain model grid
EOF
# Convert to image file using GhostScript
gmt psconvert Topography_VE.ps -A0.5c -E720 -Tj -Z
\end{lstlisting}

Prior to data visualization and mapping, the GMT scripts were written using its syntax and organized in the Xcode environment. The cocding language of GMT was used during scripting and several modules applied for plotting various cartographic elements. the area of interest was first selected and cut from the GEBCO grid using the following code: $'gmt grdcut GEBCO\_2019.nc -R286/300.5/0/12.5 -Gve\_relief.nc'.$

Following the first step, all other the substantial cartographic elements (image visualized in selected color palette, grid, clipped insert region in the globe map, coastlines, borders, rivers, barlegend, scale bar, text annotations and time stamp) were added in a script using the following GMT modules: psbasemap, grdimage, psscale, pstext, psclip, grdcontour. The technique of scripting was applied from the existing detailed descriptions \cite{Lemenkova2021b}. For example, the 'grdcontour' module was used for computing, modelling and visualizing topographic isolines using quantitative DEM data as follows: $'gmt grdcontour ve\_relief1.nc -R -J -C1000 -Wthinner,darkbrown -O -K >> \$ps'$. The purpose of the geoid visualization (Fig. 4) in relation to geophysical differentiation was to investigate how gravity values respond to the effects of topographic elevation and bathymetric depressions. Data on geoid undulations have been collected from the EGM2008 \cite{Pavlis}. 

\begin{lstlisting}[language=bash, caption=\textbf{Algorithm of GMT code used for plotting Fig. \ref{fig4} (geoid modelling)}]
#!/bin/sh
# Purpose: geoid of Venezuela
# GMT modules: gmtset, gmtdefaults, grdcut, makecpt, grdimage, psscale, grdcontour, psbasemap, gmtlogo, psconvert
# http://soliton.vm.bytemark.co.uk/pub/cpt-city/kst/tn/33_blue_red.png.index.html
# GMT set up
gmt set FORMAT_GEO_MAP=dddF \
# Overwrite defaults of GMT
gmtdefaults -D > .gmtdefaults
gmt grdconvert n00w90/w001001.adf geoid_02.grd
gdalinfo geoid_02.grd -stats
# Minimum=-70.856, Maximum=32.958, Mean=-24.768, StdDev=19.081
# Generate a color palette table from grid
gmt makecpt -C33_blue_red.cpt -T-55/25 > colors.cpt
# Generate a file
ps=Geoid_VE.ps
gmt grdimage geoid_02.grd -Ccolors.cpt -R286/300.5/0/12.5 -JM6.5i -P -Xc -I+a15+ne0.75 -K > $ps
# Add shorelines
gmt grdcontour geoid_02.grd -R -J -C1.0 -A2.0+f8p,0,black -Wthinner,dimgray -O -K >> $ps
# Add grid
gmt psbasemap -R -J \
    --MAP_FRAME_AXES=WEsN \
    --FORMAT_GEO_MAP=ddd:mm:ssF \
    -Bpx4f2a2 -Bpyg4f2a2 -Bsxg4 -Bsyg2 \
    --MAP_TITLE_OFFSET=0.8c \
    --FONT_ANNOT_PRIMARY=7p,0,black \
    --FONT_LABEL=7p,25,black \
    --FONT_TITLE=13p,25,black \
    -B+t"Geoid model (EGM-2008) of Venezuela and Trinidad and Tobago" -O -K >> $ps
# Add legend
gmt psscale -Dg286/-1.0+w16.5c/0.4c+h+o0.0/0i+ml+e -R -J -Ccolors.cpt \
    -Bg2f0.2a4+l"Color scale '33_blue_red': colour tables of IDL from the KDE scientific plotting tool [256, discrete, RGB, 252 segments, -T-55/25]" \
    -I0.2 -By+lm -O -K >> $ps
# Add scale, directional rose
gmt psbasemap -R -J \
    --FONT=7p,0,black \
    --FONT_ANNOT_PRIMARY=6p,0,black \
    --MAP_TITLE_OFFSET=0.1c \
    --MAP_ANNOT_OFFSET=0.1c \
    -Lx14.7c/-2.3c+c50+w200k+l"Mercator projection. Scale (km)"+f \
    -UBL/-5p/-65p -O -K >> $ps
# Add coastlines, borders, rivers
gmt pscoast -R -J -P -Ia/thinnest,blue -Na -N1/thick,brown -Wthick,darkslategray -Df -O -K >> $ps
# Texts
gmt pstext -R -J -N -O -K \
-F+jTL+f10p,21,darkred+jLB+a-0 >> $ps << EOF
290.0 10.5 (Coro)
EOF
# Add GMT logo
gmt logo -Dx7.0/-2.9+o0.1i/0.1i+w2c -O -K >> $ps
# Add subtitle
gmt pstext -R0/10/0/15 -JX10/10 -X0.1c -Y5.9c -N -O \
    -F+f10p,25,black+jLB >> $ps << EOF
3.0 13.6 World geoid image EGM2008 vertical datum 2.5 min resolution
EOF
# Convert to image file using GhostScript
gmt psconvert Geoid_VE.ps -A1.0c -E720 -Tj -Z
\end{lstlisting}

Whilst GEBCO data were gathered in NetCDF format ($GEBCO\_2019.nc$), the geoid data needed to be converted first from the .adf format which was done by the following code of GMT: $'gmt grdconvert n00w90/w001001.adf geoid\_02.grd'.$ The data extent and range were examined by GDAL and interpreted for the color palette as follows: $'gmt makecpt -C33_blue\_red.cpt -T-55/25 > colors.cpt'.$ The geoid was modelled using the following code: $'gmt grdimage geoid\_02.grd -Ccolors.cpt -R286/300.5/0/12.5 -JM6.5i -P -Xc -I+a15+ne0.75 -K > \$ps'.$

\begin{lstlisting}[language=bash, caption=\textbf{Algorithm of GMT code used for plotting Fig. \ref{fig5} (free-air gravity modelling)}]
#!/bin/sh
# Purpose: mapping free-air gravity anomaly of Venezuela
# GMT modules: gmtset, gmtdefaults, img2grd, makecpt, grdimage, psscale, grdcontour, psbasemap, gmtlogo, psconvert, pscoast
# http://soliton.vm.bytemark.co.uk/pub/cpt-city/pj/4/index.html
# GMT set up
gmt set FORMAT_GEO_MAP=dddF \
    MAP_FRAME_PEN=dimgray \
# Overwrite defaults of GMT
gmtdefaults -D > .gmtdefaults
gmt img2grd grav_27.1.img -R286/300.5/0/12.5 -Ggrav_VE.grd -T1 -I1 -E -S0.1 -V
gdalinfo grav_VE.grd -stats
# Minimum=-218.668, Maximum=942.457, Mean=5.694, StdDev=64.118
# Generate a color palette table from grid
gmt makecpt -Chaxby -T-200/200 > colors.cpt
# Generate a file
ps=Grav_VE.ps
gmt grdimage grav_VE.grd -Ccolors.cpt -R286/300.5/0/12.5 -JM6.0i -P -I+a15+ne0.75 -Xc -K > $ps
# Add isolines
gmt grdcontour grav_VE.grd -R -J -C50 -A100 -Wthinner -O -K >> $ps
# Add grid
gmt psbasemap -R -J \
    --MAP_FRAME_AXES=WEsN \
    --FORMAT_GEO_MAP=ddd:mm:ssF \
    -Bpx4f2a2 -Bpyg4f2a2 -Bsxg4 -Bsyg2 \
    -B+t"Free-air gravity anomaly for Venezuela, Trinidad and Tobago" -O -K >> $ps 
# Add legend
gmt psscale -Dg286/-1.0+w15.0c/0.4c+h+o0.0/0i+ml+e -R -J -Ccolors.cpt \
    --FONT_LABEL=7p,Helvetica,black \
    --FONT_ANNOT_PRIMARY=7p,0,black \
    --FONT_TITLE=8p,25,black \
    -Bg50f5a50+l"Color scale 'haxby' B. Haxby's color scheme for geoid & gravity [C=RGB -T-200/200]" \
    -I0.2 -By+l"mGal" -O -K >> $ps
# Add scale, directional rose
gmt psbasemap -R -J \
    -Lx14.0c/-2.3c+c50+w200k+l"Mercator projection. Scale (km)"+f \
    -UBL/-5p/-70p -O -K >> $ps
# Add coastlines, borders, rivers
gmt pscoast -R -J -P -Ia/thinnest,blue -Na -N1/thick,white -Wthin,darkslategray -Df -O -K >> $ps
# Texts
gmt pstext -R -J -N -O -K \
-F+jTL+f14p,32,darkgreen+jLB+a-330 >> $ps << EOF
290.8 7.8 Los Llanos
EOF
# Add GMT logo
gmt logo -Dx7.0/-2.9+o0.1i/0.1i+w2c -O -K >> $ps
# Add subtitle
gmt pstext -R0/10/0/15 -JX10/10 -X0.0c -Y4.8c -N -O \
    -F+f10p,25,black+jLB >> $ps << EOF
1.0 13.6 Global gravity grid from CryoSat-2 and Jason-1, 1 min resolution, SIO, NOAA, NGA.
EOF
# Convert to image file using GhostScript
gmt psconvert Grav_VE.ps -A0.5c -E720 -Tj -Z
\end{lstlisting}

Geophysical data of gravity grids (Figs. 5 and 6) were acquired from available sources \cite{Sandwell}. The seismic data were downloaded from the IRIS database using spatial extent of the Venezuela, Trinidad and Tobago and surroundings, respectively (Fig. 7). The gravity grids were also converted from the initial raw IMG format to the GRD format compatoble with the GMT as follows: $'gmt img2grd grav\_27.1.img -R286/300.5/0/12.5 -Ggrav\_VE.grd -T1 -I1 -E -S0.1 -V'$ (here for Fig. 5). The same procedure has been repeated for the Fig. 6 (vertical gradient of gravity) and followed the example of published work \cite{Lemenkova2021d}. The annotations have been added on all the maps using the 'pstext' module of GMT as in the following example: $'gmt pstext -R -J -N -O -K -F+f10p,0,black+jLB+a-0 >> \$ps << EOF 294.20 8.24 Ciudad Bolívar EOF'.$ 

\begin{lstlisting}[language=bash, caption=\textbf{Algorithm of GMT code used for plotting Fig. \ref{fig6} (vertical gradient of gravity modelling)}]
#!/bin/sh
# Purpose: geophysical mapping of Venezuela (vertical free-air gravity gradient)
# GMT modules: gmtset, gmtdefaults, img2grd, makecpt, grdimage, psscale, grdcontour, psbasemap, gmtlogo, psconvert, pscoast
# http://soliton.vm.bytemark.co.uk/pub/cpt-city/pj/4/index.html
# GMT set up
gmt set FORMAT_GEO_MAP=dddF \
# Overwrite defaults of GMT
gmtdefaults -D > .gmtdefaults
gmt img2grd curv_27.1.img -R286/300.5/0/12.5 -Ggrav_v_VE.grd -T1 -I1 -E -S0.1 -V
gdalinfo grav_v_VE.grd -stats
# Minimum=-488.036, Maximum=870.350, Mean=0.358, StdDev=43.643
# Generate a color palette table from grid
gmt makecpt -Chaxby -T-100/100 > colors.cpt
# Generate a file
ps=Grav_VE_v.ps
gmt grdimage grav_v_VE.grd -Ccolors.cpt -R286/300.5/0/12.5 -JM6.0i -P -I+a15+ne0.75 -Xc -K > $ps
# Add isolines
gmt grdcontour grav_v_VE.grd -R -J -C100 -Wthinnest -O -K >> $ps
# Add grid
gmt psbasemap -R -J \
    --MAP_FRAME_AXES=WEsN \
    --FORMAT_GEO_MAP=ddd:mm:ssF \
    -Bpx4f2a2 -Bpyg4f2a2 -Bsxg4 -Bsyg2 \
    -B+t"Vertical gravity gradient for Venezuela, Trinidad and Tobago" -O -K >> $ps  
# Add legend
gmt psscale -Dg286/-1.0+w15.0c/0.4c+h+o0.0/0i+ml+e -R -J -Ccolors.cpt \
    -Bg25f2.5a25+l"Color scale 'haxby' B. Haxby's color scheme for geoid & gravity [C=RGB -T-200/200]" \
    -I0.2 -By+l"mGal" -O -K >> $ps
# Add scale, directional rose
gmt psbasemap -R -J \
    -Lx14.0c/-2.3c+c50+w200k+l"Mercator projection. Scale (km)"+f \
    -UBL/-5p/-70p -O -K >> $ps
# Add coastlines, borders, rivers
gmt pscoast -R -J -P -Ia/thinnest,blue -Na -N1/thick,white -Wthin,darkslategray -Df -O -K >> $ps
# Texts
gmt pstext -R -J -N -O -K \
-F+jTL+f12p,21,white+jLB+a-315 >> $ps << EOF
286.7 5.8 Los Andes
EOF
# Add GMT logo
gmt logo -Dx7.0/-2.9+o0.1i/0.1i+w2c -O -K >> $ps
# Add subtitle
gmt pstext -R0/10/0/15 -JX10/10 -X0.0c -Y4.8c -N -O \
    -F+f10p,25,black+jLB >> $ps << EOF
1.0 13.6 Global gravity grid from CryoSat-2 and Jason-1, 1 min resolution, SIO, NOAA, NGA.
EOF
# Convert to image file using GhostScript
gmt psconvert Grav_VE_v.ps -A0.5c -E720 -Tj -Z
\end{lstlisting}

The earthquakes on the Fig. 7 (map of seismicity) have been plotted using the 'psxy' module of GMT by the following code: $'gmt psxy -R -J quakes\_VE.ngdc -Wfaint -i4,3,6,6s0.05 -h3 -Scc -Csteps.cpt -O -K >> \$ps'.$ Here the ‘i4,3,6,6s0.05’ means the numbers of column in the table ($quakes\_VE.ngdc$) that was converted from the original CSV from the IRIS. The volcanoes were visualized by the following code: $‘gmt psxy -R -J volcanoes.gmt -St0.4c -Gred -Wthinnest -O -K >> \$ps’$. The complex legend showing the magnitude of events was plotted using the following code: $‘gmt pslegend -R -J -Dx1.5/-3.0+w17.8c+o-2.0/0.1c  -F+pthin+ithinner+gwhite -O -K << FIN >> \$ps H 10 Helvetica Seismicity: earthquakes magnitude (M) from \\1.9 to 7.3 N 9 S 0.3c c 0.3c red 0.01c 0.5c M (7.1-7.3) <...> FIN’.$

\begin{lstlisting}[language=bash, caption=\textbf{Algorithm of GMT codes used for plotting Fig. \ref{fig7} (seismic mapping)}]
#!/bin/sh
# Purpose: seismicity map of Venezuela (earthquake distribution according to the IRIS database (1997-2021) of seismic events)
# GMT modules: gmtset, gmtdefaults, grdcut, makecpt, grdimage, psscale, grdcontour, psbasemap, gmtlogo, psconvert
# http://soliton.vm.bytemark.co.uk/pub/cpt-city/wkp/shadowxfox/tn/colombia.png.index.html
# GMT set up
gmt set FORMAT_GEO_MAP=dddF 
gmtdefaults -D > .gmtdefaults
# Extract a topographic subset of Venezuela and surroundings:
gmt grdcut ETOPO1_Ice_g_gmt4.grd -R286/300.5/0/12.5 -Gve_relief1.nc
gmt grdcut GEBCO_2019.nc -R286/300.5/0/12.5 -Gve_relief.nc
gdalinfo -stats ve_relief.nc
# Min=-4947.963 Max=5535.625
# Make color palette
gmt makecpt -Ccolombia.cpt > pauline.cpt
gmt makecpt -Cseis -T1.9/7.3/0.5 -Z > steps.cpt
ps=Seis_VE.ps
# Make raster image
gmt grdimage ve_relief.nc -Cpauline.cpt -R286/300.5/0/12.5 -JM6.5i -I+a15+ne0.75 -Xc -P -K > $ps
# Add legend
gmt psscale -Dg286/-1.0+w16.5c/0.4c+h+o0.0/0i+ml+e -R -J -Cpauline.cpt \
    -Bg500f50a500+l"Colormap..." \
    -I0.2 -By+lm -O -K >> $ps
# Add isolines
gmt grdcontour ve_relief1.nc -R -J -C500 -Wthinnest,brown -O -K >> $ps
# Add coastlines, borders, rivers
gmt pscoast -R -J -P \
    -Ia/thinner,blue -Na -N1/thickest,olivedrab1 -Wthin,lightcyan2 -Df -O -K >> $ps 
# Add grid
gmt psbasemap -R -J \
    --MAP_FRAME_AXES=WEsN \
    --FORMAT_GEO_MAP=ddd:mm:ssF \
    -Bpx4f2a2 -Bpyg4f2a2 -Bsxg4 -Bsyg2 \
    -B+t"Seismicity in Venezuela and Trinidad and Tobago: IRIS database (1997-2021)" -O -K >> $ps
# Add scale, directional rose
gmt psbasemap -R -J \
    -Lx14.0c/-3.6c+c50+w200k+l"Mercator projection. Scale: km"+f \
    -UBL/-5p/-110p -O -K >> $ps
# Add earthquake points
# separator in numbers of table: dot (.), not comma ! (British style)
gmt psxy -R -J quakes_VE.ngdc -Wfaint -i4,3,6,6s0.05 -h3 -Scc -Csteps.cpt -O -K >> $ps
# Add geological lines and points
gmt psxy -R -J volcanoes.gmt -St0.4c -Gred -Wthinnest -O -K >> $ps
# fabric and magnetic lineation picks fracture zones
gmt psxy -R -J GSFML_SF_FZ_KM.gmt -Wthicker,goldenrod1 -O -K >> $ps
gmt psxy -R -J GSFML_SF_FZ_RM.gmt -Wthicker,pink -O -K >> $ps
gmt psxy -R -J ridge.gmt -Sf0.5c/0.15c+l+t -Wthick,red -Gyellow -O -K >> $ps
gmt psxy -R -J ridge.gmt -Sc0.05c -Gred -Wthickest,red -O -K >> $ps
# tectonic plates
gmt psxy -R -J TP_Caribbean.txt -L -Wthickest,purple -O -K >> $ps
gmt psxy -R -J TP_Cocos.txt -L -Wthickest,purple -O -K >> $ps
gmt psxy -R -J TP_Nazca.txt -L -Wthickest,purple -O -K >> $ps
gmt psxy -R -J TP_South_Am.txt -L -Wthickest,purple -O -K >> $ps
# Texts
gmt pstext -R -J -N -O -K \
-F+f10p,17,black+jLB+a-0 >> $ps << EOF
293.20 10.58 Caracas
EOF
# Processed likewise for similar texts
gmt pslegend -R -J -Dx1.5/-3.0+w17.8c+o-2.0/0.1c \
    -F+pthin+ithinner+gwhite \
    --FONT=8p,black -O -K << FIN >> $ps
H 10 Helvetica Seismicity: earthquakes magnitude (M) from 1.9 to 7.3
N 9
S 0.3c c 0.3c red 0.01c 0.5c M (7.1-7.3)
S 0.3c c 0.3c tomato 0.01c 0.5c M (6.4-7.0)
S 0.3c c 0.3c orange 0.01c 0.5c M (5.6-6.3)
S 0.3c c 0.3c yellow 0.01c 0.5c M (4.8-5.5)
S 0.3c c 0.3c chartreuse1 0.01c 0.5c M (4.1-4.7)
S 0.3c c 0.3c chartreuse1 0.01c 0.5c M (3.4-4.0)
S 0.3c c 0.3c cyan3 0.01c 0.5c M (2.7-3.3)
S 0.3c c 0.3c blue 0.01c 0.5c M (1.9-2.6)
S 0.3c t 0.3c red 0.03c 0.5c Volcanoes
FIN
# Add GMT logo
gmt logo -Dx7.0/-4.5+o0.1i/0.1i+w2c -O -K >> $ps
# Add subtitle
gmt pstext -R0/10/0/15 -JX10/10 -X0.5c -Y5.9c -N -O \
    -F+f11p,25,black+jLB >> $ps << EOF
1.0 13.6 DEM: SRTM/GEBCO, 15 arc sec grid. Earthquakes: IRIS Seismic Event Database
EOF
# Convert to image file using GhostScript
gmt psconvert Seis_VE.ps -A1.7c -E720 -Tj -Z
\end{lstlisting}

The full algorithms of GMT codes are available in author’s GitHub repository: \\\href{https://github.com/paulinelemenkova/Mapping_Venezuela_GMT_Scripts}{https://github.com/paulinelemenkova/Mapping\_Venezuela\_GMT\_Scripts}. 

%%%%%%%%%%%%%%%%%%%%%%%%%%%%%%%%%%%%%%%%%%
\section{Results}

Fig. 1 has been plotted using GEBCO raster grid selected due to its quality and reliability. Precise topography is fundamental to the understanding of geophysics and tectonics and associated effects on ocean and terrestrial geomorphology, mapping seismicity, earthquake localization, risk assessment and tsunami wave propagation. Here the topographic and bathymetric elevations of the study area encompassing Venezuela, Trinidad and Tobago islands range from -4948 m to 5536 m according to the GEBCO grid. 

Figs. 2 and 3 indicate spatial extent of geological provinces and outcrops in Venezuela. The map in Fig. 2 shows that most of the territory of Venezuela is covered by Guyana Shield (aquamarine color) in the south-eastern parts of the country. The Barinas-Apure sedimentary basin is located in the Llanos region and the Andean foothills region in the western Venezuela \cite{Callejon}. The Barinas-Apure basin (yellow color in the map) includes oil fields that is the 3rd most important producing area in Venezuela after the Maracaibo and Eastern Venezuela basins according to the petroleum accumulation. The Falcón sedimentary basin, which produces indigenous Miocene oil from structural and stratigraphic accumulations \cite{Rodriguez} is coloured slate blue in the map. 

The South Caribbean deformed belt (purple color in Fig. 2) represents a submarine prism formed between the subducting tectonic plates in the Colombian and Venezuelan basins and arc terranes along the northern rim of South America \cite{Mann}. Its SW part, the Sinu fold belt, forms an accretionary wedge in the place of the subduction of the Caribbean Plate under South American Plate \cite{Rodriguez2021}. Other important geological provinces include the Guyana Basin which stretches along the passive margin of NE South America and the Maracaibo Basin (green, Fig. 2). The Cariaco Basin (beige color in Fig. 2) is an East-West stretching basin \cite{Schubert} situated in the Gulf of Cariaco, on the continental shelf of the Caribbean Sea, off the eastern coast (near the Barcelona, see Fig. 1). 

The Tobago Trough (chartreuse green color in Fig. 2) is a modern marine forearc basin of the Lesser Antilles arc located above the sedimentary succession at least 10 km thick and probable oceanic basement and formed by sediments from various stratigraphic sequences: Quaternary, early Pleistocene – late Miocene, Miocene and late Eocene and may also include mid-Cretaceous or older \cite{Speed}. The Quaternary sediments covering the most of the north-west of the country (dark pink color in Fig. 3) corresponds with the previous studies on Plio-Quaternary extension in the Venezuelan Andes showing the extensional structures corresponding to the elongated tilted blocks with geometry and kinematics of the structures corresponding to the earlier syn-orogenic extension \cite{Dhont}; Backé et al. 2006).

Fig. 4 show the variations of the geoid heights over the study area. The highest values of geoid (up to 25 m, bright red colors in Fig. 4) are visible in the south-western region of the map (Colombia) while the majority of Venezuela covers the extent of -50 to 8 m with the highest values in the mountains and clearly visible correlation with topographic elevations. Thus, if compared geoid mode in Fig. 4 with topographic and geologic maps in Figs 1 and 2, respectively, the elevated values of geoid are notable in the region of the Guyana Shield. The Guyana Shield is a Precambrian geological formation in the NE coast of the South America and an area of special importance in historical geology of Venezuela due to high potential for mineral exploration \cite{Gibbs}. 

\begin{figure}[H]\centering
	\includegraphics[width=12.0 cm]{Fig_04}
	\caption{Geoid model over the Venezuela, Trinidad and Tobago. Mapping: GMT. Source: author.\label{fig4}}
\end{figure}  

The environmental importance of this unique region consists in the highest biodiversity in the world including many endemic species \cite{Hollowell} in tropical forests of Amazonia \cite{Hammond}. The undulations of geoid in the Guyana Highlands corresponds to the higher topographic elevations on the shield as table-like mountains ‘tepuis’. The marine areas of the Caribbean Sea have in general lower geoid heights (blue colored areas in Fig. 4) which well corresponds to the bathymetric depressions reflected in the geophysical setting of the Earth.

Fig. 5 shows variations in the gravity over the study area. The notable free-air gravity low coincides with distribution of the Guyana basin, SE part of the Venezuela, Trinidad and Tobago, which points at the tectonic basement depression formed during Jurassic rifting, and Lake Maracaibo, an important geologic region of Venezuela. The NE coasts of the Lake Maracaibo are occupied by the Bolivar Coastal Fields, the largest oil field in South America, the 2nd oil field most prolific hydrocarbon basins in the world after the Middle East with basin area of 50,000 km2 and oil production of 30 *109 bbl \cite{Escalona}. 

The Maracaibo Lake has a wide variety of oil reservoir rocks with structural traps in the tectonic units, e.g. normal or inverted faults on the South American plate \cite{Bockmeulen}. The oil formed in the Maracaibo basin is mostly heavy black oil with asphaltene deposition \cite{Escobar,Memon}. The gravity in the Lake Maracaibo reach below -200 mGal correlating to the bathymetric depression (Fig. 1) and geologic extend of the Maracaibo Basin (Fig. 2). the highest values for gravity reach +200 mGal and well corresponds to the topographic elevation of the Cordillera de Mérida, Columbian Andes, Guiana Highlands and Sierra Parima. The Amazon and Orinoco Flow basin mostly have gentle fluctuations of gravity with values in the range between -25 to 30 mGal (Fig. 5).

\begin{figure}[H]\centering
	\includegraphics[width=12.0 cm]{Fig_05}
	\caption{Free-air gravity model of the Venezuela, Trinidad and Tobago. Mapping: GMT. Source: author.\label{fig5}}
\end{figure}   

Fig. 6 refers to a vertical gravity gradient of normal gravity as an approximation in which a gravity measurement point, which is almost never placed on the reference ellipsoid’s surface is adjusted by geophysical corrections \citep{Li,Yilmaz}. The details regarding the distribution of vertical gravity values over the Venezuela and the Trinidad and Tobago showed that the lowest values (below -100 mGal) mostly correspond to the extreme amplitude of topographic changes (e.g. abrupt border between the mountain hills and margins of the bathymetric depression) whilst the highest values (over 100 mGal) were detected in the top of the mountains, well corresponding to the local geomorphology of the Venezuela and the Trinidad and Tobago region.

\begin{figure}[H]\centering
	\includegraphics[width=12.0 cm]{Fig_06}
	\caption{Vertical gradient of gravity in the Venezuela, Trinidad and Tobago. Mapping: GMT. Source: author.\label{fig6}}
\end{figure}  

The distribution of the earthquakes and seismicity of the Venezuela and the Trinidad and Tobago (Fig. 7) depends on the proximity of area to the lithosphere plate border with active margin such as South American and Caribbean plates (depicted in the map as purple thick line). The second factor includes the geologic active processes (faults), and geomorphology. Hence, Cordillera de Mérida has the majority of the detected earthquakes following the Trinidad and Tobago active zone. The majority of mapped earthquakes I nVenezuela belong to the category ‘shallow’, that is, having depth between 0 and 70 km. Earthquakes at depth 74 to 155 km are located in the offshore Sucre, Venezuela, those detected in the Trinidad and Tobago are dominantly shallow not exceeding 70 km. The deepest earthquakes from the inspected database are located in the northern Columbia (182-200 km).

\begin{figure}[H]\centering
	\includegraphics[width=12.0 cm]{Fig_07}
	\caption{Seismic map of earthquakes in the region of Venezuela, Trinidad and Tobago. Mapping: GMT. Source: author.\label{fig7}}
\end{figure}   

\section{Conclusion}

\subsection{Summary}

This study evaluated machine learning scripting cartographic methods of GMT to assess geophysical and geologic correspondence for seismically active region of selected countries in northern Southern America – Venezuela and the Trinidad and Tobago, the Caribbean. The study applied several high-resolution geophysical, topographic and geological datasets using GEBCO, EGM-2008, remote sensing gravity grids and IRIS seismic data. The results were corresponding and supporting previous interdisciplinary studies of geological and geophysical methods of gravity measurements, active seismic observations and earthquake monitoring in Venezuela and the Caribbean Sea basin \citep{Guenther,Schmitz2008}.  

The study presented improved machine-based mapping workflow and thus contributes to the regional studies of the eastern Caribbean and Venezuela. Of the GMT models tested, several code snippets are presented as a demonstration of the technical possibilities of toolset. The results showed the best cartographic performance and demonstrate the advantage of automating processing of geophysical datasets over the GIS. The motivation to use multi-source geophysical, geologic, seismic and topographic data was to allow for comparative analysis of interplay among the processes and phenomena. The tectonic evolution of the Caribbean Sea basin largely affects its geological setting that can be tracked in geophysical data and indicate high seismicity in the Andes pointing at earthquake risk.

Since eastern Caribbean is a tectonically complex area which includes a subduction zone, island arc and active volcanoes \cite{Ambeh}, regional analysis and visualization of the varied setting was performed. Specifically, high seismicity was detected in the north-west of Venezuela, and north-east region of Trinidad and Tobago, which is consistent with geologic structure associated with tectonic setting, and related movements of the Caribbean Plate subduction. Although continued research is needed to improve earthquake monitoring in Venezuela and Trinidad and Tobago, this study presented a contribution to seismic studies in the Caribbean and South America, showing earthquakes detected according to the IRIS database based on the 25-year time span (1997-2021). 
	
Scripting cartography is currently rapidly developing and will remain the best technical approach to map seismically active regions in real time regime in order to perform prognosis, predictive mapping, risk assessment and hazard visualization. The best known advantages of scripting cartography include automatization that enables to process big datasets rapidly and effectively using multi-source heterogeneous data in massive volumes. The second advantage consists in the machine based graphics that results in aesthetically fine plotting. Besides the technical advantages of the presented scripting approach of GMT, a better understanding of geological and tectonic factors affecting seismicity and regulating the appearance, magnitude and focal depth of the earthquakes is critical to seismically active regions, such as the Caribbean and Venezuela and geologic risk assessment in South America.

\subsection{Highlights for public administration, city management and urban planning}

\begin{enumerate}
  \item The region of Venezuela and Trinidad and Tobago is at high risk of seismicity and earthquakes, as both countries are located in the zone of the South America and Caribbean tectonic plates collision and Andean orogeny. 
  \item Geological disasters from high seismicity are a function of hazard, exposure, and vulnerability of people at risk due to the location in the vicinity of earthquakes. This requires detailed studies and visualization of regional geophysical setting in the country.
  \item Integration and visualisation of the multi-source data supports geophysical investigations aimed to prevent and evaluate risks of geophysical hazards, because they provide detailed insights into the environmental, geological and topographic settings of Venezuela and supports complex spatial analysis in this seismically active region of South America. 
  \item Technically, geophysical and geological analysis of regions of seismic risk can be performed using combination of scripting Generic Mapping Tools (GMT) and Geographic Information Systems (GIS) for effective cartographic data analysis in real-time regime for hazard mapping and spatial analysis of the regions at risk.
\end{enumerate}

%%%%%%%%%%%%%%%%%%%%%%%%%%%%%%%%%%%%%%%%%%
\vspace{6pt} 

\dataavailability{Yes, the GitHub repository of the author with available GMT scripts: \href{https://github.com/paulinelemenkova/Mapping_Venezuela_GMT_Scripts}{https://github.com/paulinelemenkova/Mapping\_Venezuela\_GMT\_Scripts}} 

\acknowledgments{The author is thankful to the anonymous reviewers for reading and commenting on this study.}

\conflictsofinterest{The author declares no conflict of interest.} 

\abbreviations{Abbreviations}{
The following abbreviations are used in this manuscript:\\

\noindent 
\begin{tabular}{@{}ll}
GIS & Geographic Information System \\
GMT & Generic Mapping Tools \\
QGIS & Quantum GIS \\
GEBCO & The General Bathymetric Chart of the Oceans \\
EGM & Earth Gravitational Models \\ 
IRIS & Incorporated Research Institutions for Seismology \\
USGS & United States Geological Survey \\
CSV & comma-separated values \\
NGDC & National Geophysical Data Center
\end{tabular}}

%%%%%%%%%%%%%%%%%%%%%%%%%%%%%%%%%%%%%%%%%%


%%%%%%%%%%%%%%%%%%%%%%%%%%%%%%%%%%%s%%%%%%%
\begin{adjustwidth}{-\extralength}{0cm}
%\printendnotes[custom] % Un-comment to print a list of endnotes

\reftitle{References}

% Please provide either the correct journal abbreviation (e.g. according to the “List of Title Word Abbreviations” http://www.issn.org/services/online-services/access-to-the-ltwa/) or the full name of the journal.
% Citations and References in Supplementary files are permitted provided that they also appear in the reference list here. 

%=====================================
% References, variant A: external bibliography
%=====================================
\bibliography{Venezuela.bib}
\nocite{*}
 
%%%%%%%%%%%%%%%%%%%%%%%%%%%%%%%%%%%%%%%%%%
\end{adjustwidth}
\end{document}